\chapter{Problem Formulation \& Evaluation Metric Analysis}
\label{chap:problem}

\section{Mathematical Modeling of Reading Utility}
In this personalized recommendation setting, each candidate research paper is evaluated according to the researcher's discrete preference rating $y \in \{1, 2, 3, 4, 5\}$:
\begin{itemize}
  \item \textbf{Rating 1 (Do Not Read):} Publications falling completely outside the researcher's domain, or non-research conference metadata (such as conference proceedings headers, workshop prefaces, or organizing committee overviews). Recommending these wastes cognitive bandwidth.
  \item \textbf{Rating 2 (Low Priority):} Publications offering tangentially relevant background knowledge, suitable for occasional skimming.
  \item \textbf{Rating 3 (Medium Priority):} Papers containing standard technical methods or formalisms pertinent to the researcher's broad discipline.
  \item \textbf{Rating 4 (High Priority):} Highly relevant contributions authored by prominent peers or addressing key theoretical challenges directly aligned with the researcher's agenda.
  \item \textbf{Rating 5 (Must Read):} Groundbreaking papers providing critical extensions or direct solutions to the researcher's core topics; must be read immediately.
\end{itemize}

The recommendation system must assign a continuous priority score $\text{priority}_i \in [0, +\infty)$ to each candidate paper in the unrated test collection:
\begin{itemize}
  \item $\text{priority}_i = 0$: Exclude paper $i$ from the recommendation list.
  \item $\text{priority}_i > 0$: Recommend paper $i$, with strictly larger numerical values dictating earlier reading positions.
\end{itemize}

\section{Formal Evaluation Metric Breakdown}
The 596 test publications are partitioned into two equal, disjoint hidden splits: the \textbf{Public Test set (298 papers)} and the \textbf{Private Test set (298 papers)}. Within each split independently, the competition evaluator executes the following procedure:
\begin{enumerate}
  \item Excludes all candidate papers assigned $\text{priority} \le 0$.
  \item Sorts the remaining $m$ positive-priority papers in descending order of priority, breaking numerical ties by ascending alphabetical paper ID.
  \item Assigns sequential reading rank positions $r = 1, 2, \dots, m$ and computes the official competition score.
\end{enumerate}

The final score is the Harmonic Mean of Relevance ($Q$) and Coverage ($R$):
\begin{equation}
  \text{Score} = \frac{2 \cdot Q \cdot R}{Q + R}
  \label{eq:harmonic_mean}
\end{equation}

\subsection{Relevance ($Q$): Discounted Cumulative Gain}
Relevance measures the positional utility of the recommended reading sequence. Following the Discounted Cumulative Gain (DCG) formulation \cite{jarvelin2002cumulated}, earlier positions receive substantially greater weight:
\begin{equation}
  Q = \frac{\sum_{r=1}^m w_r \cdot g_r}{\sum_{r=1}^m w_r}
  \label{eq:relevance}
\end{equation}
where:
\begin{itemize}
  \item $g_r = \frac{y_r - 1}{4} \in [0, 1]$ represents the normalized usefulness gain for a paper with hidden rating $y_r$. The exact gain mapping is summarized in Table \ref{tab:gain_weights}.
  \item $w_r = \frac{1}{\log_2(r + 1)}$ denotes the position discount weight.
\end{itemize}

\begin{table}[!htbp]
  \centering
  \small
  \begin{tabularx}{0.85\linewidth}{>{\centering\arraybackslash}p{2.5cm} >{\centering\arraybackslash}p{2.5cm} >{\raggedright\arraybackslash}X}
    \toprule
    \rowcolor{PrismHeader}
    \color{PrismWhite}\textbf{True Rating ($y_r$)} & 
    \color{PrismWhite}\textbf{Usefulness Gain ($g_r$)} & 
    \color{PrismWhite}\textbf{Contribution to Relevance ($Q$)} \\
    \midrule
    1 & $0.00$ & Contributes zero gain, inflates denominator $\to$ \textbf{Severely penalizes $Q$} \\
    \rowcolor{PrismLabel}
    2 & $0.25$ & Minimum positive gain ($25\%$) \\
    3 & $0.50$ & Moderate utility gain ($50\%$) \\
    \rowcolor{PrismLabel}
    4 & $0.75$ & High utility gain ($75\%$) \\
    5 & $1.00$ & Maximum possible gain ($100\%$) \\
    \bottomrule
  \end{tabularx}
  \caption{Mapping from true reading preference ratings to normalized usefulness gain $g_r$}
  \label{tab:gain_weights}
\end{table}

\begin{prismbox}[title=\textbf{The Severe Mathematical Penalty of Rating-1 Inclusions}]
Consider what happens when a rating-1 paper ($y_r = 1 \implies g_r = 0$) infiltrates the recommendation list at position $r$:
\begin{itemize}
  \item The numerator $\sum w_r g_r$ increases by $w_r \times 0 = 0$.
  \item The denominator $\sum w_r$ strictly increases by $w_r > 0$.
\end{itemize}
Consequently, the fraction drops immediately. Because $w_1 = 1.0, w_2 \approx 0.63$, and $w_3 = 0.50$, allowing a rating-1 paper into the top ranks acts as a catastrophic drag on $Q$. Therefore, \textbf{the foremost architectural priority of the system is the decisive suppression of all suspected rating-1 items}.
\end{prismbox}

\subsection{Coverage ($R$): Recall of Worthwhile Publications}
Coverage quantifies the recall of worthwhile papers (ratings $y \ge 2$) retained in the recommendation pool:
\begin{equation}
  R = \frac{\sum_{r=1}^m \mathbb{I}(y_r > 1)}{N_{\text{worthwhile}}}
  \label{eq:coverage}
\end{equation}
where $N_{\text{worthwhile}} = \sum_{j=1}^N \mathbb{I}(y_j > 1)$ represents the total count of worthwhile papers in the evaluated split, and $\mathbb{I}(\cdot)$ is the indicator function. Key characteristics of Coverage include:
\begin{itemize}
  \item Rank invariance: Reordering selected papers does not alter $R$.
  \item Zero reward for rating 1: Recommending rating-1 papers adds nothing to the numerator.
  \item False negative penalty: Erroneously filtering out a rating $\ge 2$ paper directly reduces $R$.
\end{itemize}

\subsection{Harmonic Mean \& The Fundamental Precision-Recall Trade-off}
The selection of the Harmonic Mean over the Arithmetic Mean creates an acute trade-off:
\begin{itemize}
  \item \textbf{High Filtering Threshold $T$ (Aggressive Pruning):} Extremely high precision eliminates nearly all rating-1 items, driving $Q$ above $0.65$. However, marginal rating-2 and rating-3 papers are dropped, crashing $R$ below $0.50$ and depressing the final Harmonic Mean.
  \item \textbf{Low Filtering Threshold $T$ (Lenient Inclusion):} High recall secures $R > 0.90$, but rating-1 papers inevitably leak into the ranking list, diluting $Q$ below $0.45$ and causing a severe overall score drop.
\end{itemize}
Hence, optimal performance hinges on tuning a precise decision threshold $T^*$ that maximizes the harmonic balance on out-of-fold predictions.

\section{Theoretical Ceiling Analysis}
To establish the maximum performance achievable under perfect information, we executed a Monte Carlo simulation (100 trials) on 298-paper splits using a hypothetical \textit{Oracle Recommender} with omniscient knowledge of ground truth labels:
\begin{itemize}
  \item The Oracle filters all rating-1 papers by assigning $\text{priority} = 0$.
  \item It places all rating-5 papers at ranks $1, \dots, N_5$, followed by rating-4, rating-3, and rating-2 papers.
\end{itemize}

The simulation establishes an absolute upper bound:
\begin{equation}
  \text{Score}_{\text{Oracle}} \approx \mathbf{0.77654} \quad (Q = 0.63489, \; R = 1.00000)
\end{equation}

\begin{prismbox}[title=\textbf{Why Is the Maximum Theoretical Ceiling 0.776 Rather Than 0.800?}]
For the final score to reach $\text{Score} \ge 0.800$ with $R = 1.000$:
\begin{equation*}
  \frac{2Q}{Q + 1} \ge 0.8 \iff 1.2Q \ge 0.8 \iff Q \ge \frac{2}{3} \approx \mathbf{0.6667}
\end{equation*}
However, within the actual empirical distribution, rating-2 ($g=0.25$) and rating-3 ($g=0.50$) papers constitute nearly $\mathbf{60\%}$ of all worthwhile publications. Because securing $R = 1.0$ mandates their inclusion, their modest gain values inevitably dilute the position-weighted average $Q$, capping it at $\mathbf{0.6349}$.

If the Oracle attempts to artificially inflate $Q$ to $0.768$ by dropping all rating-2 papers, Coverage drops precipitously to $R = 0.674$, causing the overall score to drop from $0.776$ down to $\mathbf{0.717}$. Thus, $\mathbf{0.776}$ represents the empirical theoretical ceiling of this dataset.
\end{prismbox}
